\section{Composite degrees and the contributions of multiple images}
\label{sec:composite}
Write $p_{h,d}^\epsilon=[q^d]\log F_h^\epsilon(q)$ for a universal
connected class and $S_{h,d}^\epsilon$ for its numerical connected
residue, $\epsilon=\mathrm{sc},\mathrm{ad}$. The power-locus bound in
Theorem~\ref{thm:main} excludes cycles that are proper integral multiples of another entire cycle. A cycle can still have an individually repeated component. We now retain the double of an entire cycle and determine its contribution.
In the formulas below $j$ may have either sign.

\begin{theorem}\label{thm:doubleimage}
Suppose $m\ge2$, $|j|\le n_m$ and $|j|>n_{m/e}$ for every divisor
$e\ge3$ of $m$. For even $m$ write
$P_{m/2,j}=\sum_h b_h x_{h,1}$, interpreting a weight outside the
Laurent range as zero. For odd $m$ set $b_h=0$. Then there is a
finite polynomial $R$ such that
\[
 P_{m,j}=\sum_h b_hp_{h,2}^{\rm sc}+\iota(R),\qquad
 R=\sum_kN_{k,m,j}u^{2k-2}-\sum_h b_h S_{h,2}^{\rm sc}.
\]
The coefficient of $w^{h-1}$ in $R$ vanishes if
$h<1+m(g-1)+|j|$. In particular this computes every nonconstant
quartic weight in every core genus and the sextic weights $|j|>g-1$.
\end{theorem}

\begin{proof}
On the projective base for the $j$th coefficient the translated
spectral tests avoid $e$th powers for each $e\ge3$ by the
codimension computation of Section~\ref{sec:weights}. Avoidance
persists in a collar of the compact incidence. If $m$ is even, all
proper multiples in degree $m/2$ are similarly excluded, so
Theorem~\ref{thm:main} supplies the $b_h$ when $m/2>1$.
In degree one the only weight is zero. Proper division by two
sends $P_{m,j}$ to $P_{m/2,j}$ and $p_{h,2}^{\rm sc}$ to $x_{h,1}$.
All remaining proper divisions vanish on both sides. The logarithms
are primitive. It follows from multiplicity-one detection that the
remainder has the form $\iota(R)$ for a unique finite Laurent
polynomial with exponents at least $-1$.

Fix a bound $K$ on source genus. In the separation argument of
Theorem~\ref{thm:separate} keep every tuple of distinct components
with normalization genus sum at most $K$, including the root of a
double whose own degree-two covers have minimum genus above $K$.
At each step order by total reduced degree, the genus sum and then
minus the number of distinct images. Lemma~\ref{lem:weighted-specialization}
makes the first remaining stratum compact after the neighborhoods
already treated have been excised. The shear and jet argument makes
its normalized images embedded and pairwise disjoint. A connected
support is therefore a simple embedded curve or a double of one.
The equations on a reduced Hilbert chart are those of its
\emph{weighted} cycle. Its transverse solutions are finite.

Before excising a doubled solution, fix its reduced curve $D'$ and
vary the normal bundle. Here $\det N_{D'}=K_{D'}$, after trivializing
a constant line. At an unramified fiber away from the tests choose
a framed chart, a uniformizer $t$, a linear form $\ell$ separating
all sheets $y_i$, and the monic polynomial
$S(T)=\operatorname{Norm}(T-\ell(y))$ of their union on the chart.
For a constant area bivector choose $v=\iota_\ell\Omega$, so
$\ell(v)=0$. The transition
\[
 y\longmapsto y+v S(\ell(y))
          \sum_{a=1}^{B}\sum_{l=0}^{r-1}c_{al}t^{-a}\ell(y)^l
\]
is an automorphism of the punctured chart with determinant one.
It fixes each sheet. Its first normal derivative on sheet $i$ is
given by $v\otimes\ell$ times $S'(\ell(y_i))$ times the displayed
polynomial. The Vandermonde determinant is a unit modulo $t^B$ and
three generic forms span $\mathfrak{sl}_2$ in this way.
For large $B$ these principal parts map onto
$H^1(D',\operatorname{End}_0 N_{D'})$; using the union of disjoint
sheet values makes the jets independent for several marked curves.
The same framed parameter bundles and base cutoffs as in
Lemma~\ref{lem:shear-jets} realize these jets in one algebraic
target family, with a smooth section over the compact parameter base.

A good normal bundle is accessible even from an unstable one.
For an ample enough line $Q$, the initial bundle and any prescribed
good bundle of this determinant both admit quotients $Q$ (a general
section of the globally generated $N^*Q$ has no zeros). They belong
to one irreducible extension space
$\operatorname{Ext}^1(Q,K_{D'}Q^{-1})$. The open set of stable
two-rigid bundles there is dense by \cite[Theorem~2.2]{BP06}.
The preceding Kodaira--Spencer map is onto and deformations of
fixed-determinant bundles on a curve have no second obstructions.
Thus the parameter family is smooth over the bundle stack locally
and can realize an arbitrarily small change into this open.
The regular cut and its sign persist. Each root then has acyclic
normal bundle and extends uniquely on a relative Hilbert chart.
Any other nearby cycle of multiplicity at most two, recentered
and dilated to unit normal height, would limit to a compact
off-section cycle in the normal bundle. The disk mass bound follows
from the fixed projected class and a common closed taming form.
The normalization of a positive component would supply a nonzero
normal section on a cover of degree at most two. Two-rigidity excludes
this case. This proves isolation on the full cycle space, uniformly
on a smaller parameter neighborhood. The separation induction can
now proceed with these compact neighborhoods excised.

Over the isolated chart the doubled point is the scalar class with
empty evaluation tuple. Its connected pullback tests give those of
$p_{h,2}^{\rm sc}$: the product in the logarithmic subtraction
vanishes on connected domains. The local intersection degree for the
divided insertion is the same sign, since doubling maps the reduced
section homeomorphically to the doubled section. Testing the divided
insertion successively in genera through $K$ identifies its normalized
signed totals with $b_h$. The same triangular test on the residual
simple intersections identifies their totals with $[w^{h-1}]R$
through $K$.

For the virtual pairing the difference of tangent complexes at a
map $f$ to the moving reduced curve is
$\operatorname{Cone}(q^*Rp_*N\to R\pi_*f^*N)$, with the notation of
Proposition~\ref{prop:moving-embedded-curve}. The first term vanishes
by acyclicity. Two-rigidity and cohomology and base change make the
second a bundle of rank $2(k-2h+1)$ in degree one on the whole
proper stack of degree-two maps, including its nodal boundary.
Its top Chern class pushes to the scalar degree-two connected residue
times the base class: on a single fiber the equality follows by
equivariant localization, where only the top Chern number survives
in virtual dimension zero. Simple intersections contribute
$w^{h-1}$ by the same proposition. The fixed-genus pairings thus
give the asserted equality for arbitrary $K$. A nonzero double
coefficient has $h\ge1+(m/2)(g-1)+|j|$; its covers start at
$1+m(g-1)+2|j|$. Together with the spectral genus bound for the
original cut, this proves the support assertion.
\end{proof}

\subsection{The exceptional sextic coefficient}
For a genus-two core $n_6=15$. At $j=1$ the compact base is
$\PP^{16}\times\PP^{14}$. A cube of a trace-free quadratic
$T^2+b$ has coefficients in orders three through six equal to
$0,a_2^2/3,0,a_2^3/27$ where $a_2=3b$. The independent sections
enforcing these four conditions on the second factor have ranks
$2,3,4,5$, total $14$, and positive top Chern number. They cut
transversely for general translations and avoid the further condition
$a_2=0$. For general $M$ the sole divisor of $M^2$ has two distinct
points, so each reduced quadratic is smooth of genus four. Properness
of multiplication of cycles and openness near a smooth embedded
cycle show that throughout a small collar every pure triple in the
cut has such a root. Pure sixth powers are absent.

Repeat the preceding proof, retaining doubles and triples separately.
At a tripled root, first vary the normal bundle to the general stable
locus on the genus-four curve. Theorem~\ref{thm:isolated} shows that
288 cyclic families remain in the normal bundle. At a fresh unramified
fiber the same interpolation with $S(\ell(y))^2$ fixes the first
normal jets while prescribing quadratic terms on every sheet. In
particular it realizes the determinant-one shears of
Proposition~\ref{prop:isolation} simultaneously at the finitely many
curves. A general quadratic change now isolates each tripled root
while preserving the reduced transverse cut. These changes precede
excision. The full relative germ near the multiple is its moving
section by the uniform isolation in that proposition. Thus no
remaining connected support through the chosen bound is lost.

The two proper divisions give
\[
 \rho_2P_{6,1}=30x_{5,1}-5x_{6,1}+\tfrac49x_{7,1},\qquad
 \rho_3P_{6,1}=\tfrac12 x_{4,1}.
\]
Indeed, subtract $P_{3,3}=x_{7,1}/18$ from one quarter of the
genus-two cubic row, and apply the quadratic row in
\eqref{eq:arith-two}. Let $b_5=30,b_6=-5,b_7=4/9$ and put the other
$b_h$ equal to zero. The two divided tests measure the root signs
separately without automorphism factors. Fix $K\ge13$ and apply the
connected pullback tests up to genus $K$. The undivided test at
each tripled root uses $J_3=x_{4,3}+96x_{13,1}$, and its
fixed-target test has the same extra $96w^{12}$. Consequently
\[
 P_{6,1}=\tfrac12p_{4,3}^{\rm sc}+\sum_h b_hp_{h,2}^{\rm sc}
           +\iota\!\left( [z]S_{2,6}(z)-\tfrac12 S_{4,3}^{\rm sc}
                                      -\sum_h b_hS_{h,2}^{\rm sc}\right).
\]
Here $[z]$ means Laurent weight one; the remainder in parentheses
includes $48w^{12}$ in addition to the normalized simple-image counts.
This use of the nonlinear germ is essential: tests on a moving
reduced chart need not isolate the cubic amplitude in the linear
bundle, even when they do detect the amplitude at a fixed base point.
Indeed the derivative in the moving base coordinates at the root is
invertible. The implicit function theorem solves the test equations
as the amplitude moves.

\subsection{Conversion to opposite-weight factors}\label{sec:rows}
Define $t_{g,d,h}$ by
$F_g^{\rm sc}(q)=\prod_{d,h}F_h^{\rm ad}(q^d)^{t_{g,d,h}}$.
Let $D_{g,d}=S_{g,d}^{\rm sc}-S_{g,d}^{\rm ad}$. Subtracting the
proper divisors in the logarithm gives
\begin{align*}
 \sum_h t_{g,4,h}w^{h-1}
   &=D_{g,4}-\sum_h t_{g,2,h}S_{h,2}^{\rm ad},\\
 \sum_h t_{2,6,h}w^{h-1}
   &=D_{2,6}-2S_{4,3}^{\rm ad}
              -120S_{5,2}^{\rm ad}+20S_{6,2}^{\rm ad}-2S_{7,2}^{\rm ad}.
\end{align*}
The quartic formula follows since the sum over the doubled images
in Theorem~\ref{thm:doubleimage} is the entire quadratic change.
For the sextic the positive odd weights three and one contribute
$1/18$ and $4/9$, respectively, to the genus-seven double, giving
$4(1/18+4/9)=2$. The other doubles and the triple come from weight
one alone. The remaining proper factors convert scalar lower
logarithms into opposite-weight logarithms by Theorem~\ref{thm:main}.

For a genus-two core the quartic row in genera six through eleven is
\[
     (6860,-1470,1302,-317,40,-2).
\]
Its first positive Laurent coefficient starts at $1715w^5$.
The sextic row in genera eight through twenty-two is
\[
\begin{array}{c|rrrrr}
 h&8&9&10&11&12\\ \hline
 t_{2,6,h}&34606024&61350344&125323022&-36417600&62126672\\[2pt]
 h&13&14&15&16&17\\ \hline
 t_{2,6,h}&-41880636&22703802&-9395788&2975620&-715462\\[2pt]
 h&18&19&20&21&22\\ \hline
 t_{2,6,h}&128364&-16652&1476&-80&2
\end{array}
\]
For reproducibility, in each degree $d$ form the disconnected
partition trace $Z_d$ of \cite[Sections~4--8]{BP} and recover
$S_d=Z_d-d^{-1}\sum_{i=1}^{d-1}iS_iZ_{d-i}$. With $r=e^{iu}$ this
calculation takes place over $\QQ(r)$. Multiply the transposition
matrix by the unit and its iterates, express the other partition
operators in the resulting cyclic basis, and form the handle and
inverse caps as in Appendix~\ref{app:handle}. At $z=-1$ use
the hook sum, where the transposition operator need not be cyclic.
For odd weights $1,3,\ldots,2l-1$ the columns
$z^j+z^{-j}$ are $y=z+z^{-1}$ times polynomials of successive
degrees in $y^2$. Interpolate $S_d(z)-S_d(-z)$ at $z=1,\ldots,l$,
with $l=3$ for the quartic and $l=8$ for the sextic. An additional
pair of weights verifies each row. Substitution $w=2-r-r^{-1}$
gives the displayed polynomials. For instance the sole nonintegral
coefficient of $D_{2,6}$ is $376348226/3$ at $w^9$. The cubic
subtraction there is $563480/3$ and the sum of the double terms
to be subtracted is $-61440$, giving $125323022$.

In a compact class $m\alpha$ with $\alpha$ primitive in the free
quotient, substitution of the factors gives
\[
 c_{m\alpha,h}-e_{m\alpha,h}
 =\sum_{d\mid m,\ d>1}\ \sum_{d\gamma=m\alpha}\sum_g
       e_{\gamma,g}t_{g,d,h}
\]
over effective integral roots. For example, at $(m,h)=(6,8)$ one gets
\[
192\sum_{2\delta=6\alpha}c_{\delta,4}
 +14448\sum_{3\delta=6\alpha}c_{\delta,3}
 +34606024\sum_{6\gamma=6\alpha}c_{\gamma,2}.
\]
Indeed the further balanced double of a cubic change starts at
$w^8$ and the further triple of a quadratic change starts at $w^9$.
A genus at least three at a sixth root contributes no term of degree
seven in $w$, by Riemann--Hurwitz. For $m=4$ the corresponding identity at output genus six is
$24\sum_{2\delta=4\alpha}c_{\delta,3}
+6860\sum_{4\gamma=4\alpha}c_{\gamma,2}$. Here
$J_2=24w^5-2w^6$ and the further double of a genus-four correction
starts at $w^6$. No formula here identifies the inner modes on
arbitrary higher power loci.
